\chapter{Monte carlo and boh}

 [...]

\section{The Ising Model}

lattice, spins, hamiltonian, external coupling term, internal coupling term (interaction), energy configuations\TODO{arrows on a lattice}

The partition function is defined by
\[
    \dots
\]

We are now interested to study strongly correlated systems, where the interaction term is not negligible. In particular, are interested in the \textbf{critical behavior} of the system, which is characterized by a phase transition (second order) at a critical temperature $T_c$. The critical behavior is characterized by the divergence of the correlation length $\xi$ and the susceptibility $\chi$ as $T \to T_c$. We are not interested in first order phase transitions, which are characterized by a discontinuity in the first derivative of the free energy (latent heat) and a finite correlation length at the transition. We will be interested in second and infinite order phase transitions.

\begin{remark}[Lattice QCD]
    The natural way to regularize QFT is to put it on a lattice, which is the basis of lattice QCD. The Ising model is a toy model for lattice QCD, where the spins are replaced by quarks and the interaction term is replaced by the gauge field. The Ising model is a classical model, while lattice QCD is a quantum model. The lattice naturally introduces a cutoff, which is the lattice spacing. The continuum limit is obtained by taking the lattice spacing to zero, which corresponds to taking the cutoff to infinity.
\end{remark}

There is spontaneous symmetry breaking in the Ising model, which is characterized by a non-zero magnetization below the critical temperature. The hamiltonian is invariant under a global $\mathbb{Z}_2$ symmetry (if we neglect the external field), which is spontaneously broken below the critical temperature. The order parameter is the magnetization, which is zero above the critical temperature and non-zero below the critical temperature, where the system chooses one of the two degenerate ground states. The critical exponents characterize the behavior of the order parameter and the correlation length near the critical temperature.

There is no local order parameter in the Ising model, only global information about the system.\footnote{An order parameter is an observable that is zero in one phase and non-zero in another phase, which characterizes the phase transition. A local order parameter is an order parameter that is defined at each point in space, while a global order parameter is an order parameter that is defined for the whole system. The magnetization is a global order parameter, and if we have to build a local one we have to seek for a finite supported operator (it does not need information from the whole system to be defined) that is non-zero in one phase and zero in another phase.}

Using a topological description of the system, we can arrive to a lattice gauge theory, which is a quantum model. The Ising model is a classical model, while the lattice gauge theory is a quantum model. The lattice gauge theory is a generalization of the Ising model, where the spins are replaced by qubits and the interaction term is replaced by a gauge field. (Kitaev's toric code is a lattice gauge theory, which is a quantum model that exhibits topological order.) The topological order is characterized by a ground state degeneracy that depends on the topology of the system, and by the presence of anyons, which are quasiparticles that obey fractional statistics. The topological order is a global property of the system, which is not captured by local order parameters.

Integer quantum Hall effect, topological response of the system: robust to perturbations... topological information is preserved under continuous deformations of the system. The topological information is encoded in the ground state wavefunction, which is a global property of the system. The topological information is robust to perturbations, which is a consequence of the gap in the spectrum of the system. Topological quantum computing is a way to encode information in the topological properties of the system, which is robust to perturbations.

We are going tostudy the Ising model in different regimes, theoretically and numerically. after studying the classical Ising model, we will study its quantized version, and in order to do that we will need to use and understand different methods (mean field theory, renormalization group, Monte Carlo simulations, etc.) to study the Ising model in different regimes. We will also study the Ising model in different dimensions, and we will see that the critical behavior is different in different dimensions. We will study the duality of the Ising model, which is a powerful tool to study the critical behavior of the system.

\subsection{1D Ising Model}

We are going to study the free energy and its discontinuities, given by the partition function, which is defined in \TODO{add label}.

So the free energy is given by
\[
    F = - \lim_{N \to \infty} \frac{k_B T}{N} \log \mathcal{Z},
\]
where \(k_B = 1.38 \times 10^{-23} \, \text{J/K}\) is the Boltzmann constant, \(T\) is the temperature, \(N\) is the number of spins, and \(\mathcal{Z}\) is the partition function. The free energy is a thermodynamic potential that characterizes the equilibrium properties of the system. The discontinuities in the free energy correspond to phase transitions, which are characterized by a change in the order parameter.

\subsubsection{Transfer Matrix Method}

We can rewrite the partition function in terms of a transfer matrix, which is a matrix that encodes the interactions between neighboring spins:
\[
    \begin{aligned}
        \mathcal{Z} & = \sum_{\sigma_N = \pm 1} \sum_{\sigma_{N-1} = \pm 1} \dots \sum_{\sigma_1 = \pm 1} e^{-\beta H(\sigma_1, \sigma_2, \dots, \sigma_N)}                                       \\
                    & = \sum_{\sigma_N = \pm 1} \sum_{\sigma_{N-1} = \pm 1} \dots \sum_{\sigma_1 = \pm 1} T_{\sigma_N, \sigma_{N-1}} T_{\sigma_{N-1}, \sigma_{N-2}} \dots T_{\sigma_2, \sigma_1},
    \end{aligned}
\]
where \(T\) is the \(2\times 2\) transfer matrix, which is defined by
\[
    T_{\sigma_i, \sigma_{i-1}} = e^{-\beta H(\sigma_i, \sigma_{i-1})}, \quad \begin{pmatrix}
        e^{\beta J}   & e^{- \beta J} \\
        e^{- \beta J} & e^{\beta J}
    \end{pmatrix}.
\]
Using this formalism, we can write the partition function as a product of transfer matrices and then take the trace to obtain
\[
    \mathcal{Z} = \text{Tr}(T^N).
\]
So the whole problem reduces to finding the eigenvalues of the transfer matrix \(T\). The eigenvalues of \(T\) are given by its diagonalization, which is a standard linear algebra problem. We can compute them as follows:
\[
    \begin{aligned}
        \det(T - \lambda I) = 0, \\
        (e^{\beta J} - \lambda)^2 - (e^{- \beta J})^2 = 0,
    \end{aligned}
\]
which implies that the eigenvalues are
\[
    \begin{cases}
        \lambda_+ = e^{\beta J} + e^{- \beta J} = 2 \cosh(\beta J), \\
        \lambda_- = e^{\beta J} - e^{- \beta J} = 2 \sinh(\beta J).
    \end{cases}
\]
So in the end we have that the partition function is given by
\[
    \mathcal{Z} = \lambda_+^N + \lambda_-^N,
\]
such that the free energy is given by
\[
    F = - k_B T \lim_{N \to \infty} \frac{1}{N} \log(\lambda_+^N + \lambda_-^N) = - k_B T \log \lambda_+ = - k_B T \log(2 \cosh(\beta J)),
\]
since the term \(\lambda_-^N\) is negligible in the thermodynamic limit \(N \to \infty\). The free energy is a smooth function of the temperature, in the form of a logarithm of a positive and continuous function, so it is continuous and differentiable. Therefore, there is no phase transition in the 1D Ising model, which is consistent with the fact that there is no spontaneous symmetry breaking in one dimension.

\subsubsection{Correlation Functions}

If we consider the connected correlation function between spins, we can compute it using the transfer matrix formalism. A simple computation shows that
\[
    \langle \sigma_1 \sigma_N \rangle = \left(\frac{\lambda_-}{\lambda_+}\right)^{N} = \left(\tanh(\beta J)\right)^{N} \sim e^{-N/\xi},
\]
where \(\xi\) is the correlation length, which is given by
\[
    \xi = -\frac{1}{\log(\tanh(\beta J))}.
\]
The correlation length is finite for any finite temperature, which implies that there is no long-range order in the 1D Ising model. The correlation length diverges only in the limit \(T \to 0\), which corresponds to the ground state of the system. Therefore, there is no phase transition in the 1D Ising model, which is consistent with the fact that there is no spontaneous symmetry breaking in one dimension.

So in the 1D Ising model, the correlation length is finite for any finite temperature: the system does not exhibit a phase transition, and there it remains in a paramagnetic phase for any finite temperature. The system only exhibits a phase transition in the limit \(T \to 0\), where the correlation length diverges and the system becomes ordered. \TODO{add plot of correlation length vs temperature}
\TODO{add plot of correlation function vs distance for different temperatures (log log)}

\subsection{2D Ising Model}

Different definitions of the 2D Ising model: square lattice, triangular lattice, Kagome lattice, etc.\TODO{image of ising on different lattices} The 2D Ising model exhibits a phase transition at a finite temperature \(T_c\), which is characterized by a divergence of the correlation length and the susceptibility. But does the lattice matter? If we imagine the spacing \(J\) of the lattices to be the same, then the critical temperature \(T_c\) is different for different lattices, but we wont go into that.

Instead, let us imagine the triangular lattice, given two different spins in the bottom angles of the triangle, we can see that the spin in the top angle is \textit{frustrated}, since it cannot satisfy both interactions with the bottom spins. This is an example of \textbf{frustration}, which is a phenomenon that occurs in systems with competing interactions: the system cannot satisfy all interactions simultaneously, which leads to a degenerate ground state and a complex energy landscape. Frustration can lead to interesting phenomena, such as spin glasses, which are disordered magnetic systems with a complex energy landscape and slow dynamics. For now we are sticking to the square lattice, which is the simplest case, but we will see that on a triangular lattice the connection with QFT is more direct...

\subsubsection{Exact Onsager Solution}

The 2D Ising model on a square lattice was solved exactly by Lars Onsager in 1944, using the transfer matrix method and the theory of Toeplitz determinants. The exact solution shows that the 2D Ising model exhibits a phase transition at a finite temperature \(T_c\), which is characterized by a divergence of the correlation length and the susceptibility.\TODO{add drawing of correlation length vs temperature} It is a second order phase transition, which is characterized by a continuous change in the order parameter and a divergence of the correlation length.

Below the critical temperature, the system exhibits spontaneous symmetry breaking, which is characterized by a non-zero magnetization. The order parameter is the magnetization, which is zero above the critical temperature and non-zero below the critical temperature, where the system chooses one of the two degenerate ground states.

\subsubsection{Peierl's Argument}

Imagine to put a domain wall in the system, which is a line that separates two regions with different magnetization: in 1D imagine a line that separates two regions with different magnetization, one of all \(\uparrow\) and the other of all \(\downarrow\). The energy cost of creating a domain wall is proportional to its length, which is given by \(E \sim J \ell\), where \(\ell\) is the length of the domain wall. The entropy associated with the domain wall is proportional to its length, which is given by \(S \sim k_B \ell\). So the domain wall gives a contribution to the free energy and one to the entropy:\TODO{wrong computations, check}
\[
    \Delta F = \Delta E - T \Delta S \sim J \ell - k_B T \ell = \ell (J - k_B T).
\]
As we take the TD limit in 1D we will have
\[
    \Delta F \sim \ell (J - k_B T) \to -\infty,
\]
which implies that the domain wall is favorable and the system will be always disordered: no phase transition and always in a paramagnetic phase.

If we try to promote the argument to 2D, we will have an energy cost of creating a domain wall that is proportional to its length, which is given by \(E \sim J \ell\), where \(\ell\) is the length of the domain wall. The entropy associated with the domain wall is proportional to its area, which is given by \(S \sim k_B \ell^2\). So the domain wall gives a contribution to the free energy and one to the entropy:
\[
    \Delta F = \Delta E - T \Delta S \sim J \ell - k_B T \ell^2 = \ell (J - k_B T \ell).
\]
As we take the TD limit in 2D we will have a region in which the domain wall is not favorable, which implies that the system can be ordered and exhibit a phase transition at a finite temperature \(T_c\). The Peierl's argument shows that the 2D Ising model exhibits a phase transition at a finite temperature \(T_c\), which is characterized by a divergence of the correlation length and the susceptibility. It is a second order phase transition, which is characterized by a continuous change in the order parameter and a divergence of the correlation length.\TODO{check computations of the whole subsubsection, add drawing of domain wall in 1D and 2D}







\section{Motivation for Monte Carlo Methods}

Speaking about the ising model, we have seen that the 1D case is exactly solvable, while the 2D case is exactly solvable only in the absence of an external field. In general, the Ising model is not exactly solvable in higher dimensions or in the presence of an external field. Therefore, we need to use numerical methods to study the Ising model in these cases.

Even if the ising model is ideal and simple, its universal behavior is shared by many other systems, and therefore it is a good model to study the critical behavior of strongly correlated systems. The famous bootstrap program was developed with the aim of solving the 3D Ising model, which is not exactly solvable.

In the ising model, studying its partition function we can notice that the number of configurations grows exponentially with the number of spins
\[
    \mathcal{Z} = \sum_{\{\sigma_i\}} e^{-\beta H(\{\sigma_i\})} \sim 2^N,
\]
where \(N\) is the number of spins. Even if we don't consider the simulation of the system but only the computation of the partition function, we can see that the number of configurations reaches astronomical numbers for a relatively small number of spins. For example, for \(N = 100\) spins, the number of configurations is \(2^{100} \sim 10^{30}\), which is larger than the number of atoms in the observable universe, orders of magnitude more data than what we can store in a computer. Therefore, we need to use numerical methods to study the Ising model in these cases.

Speaking about \(P \neq NP\) problems, we can see that the Ising model is a \(NP\) problem, since we can verify a solution in polynomial time, but we cannot find a solution in polynomial time. The Ising model is a combinatorial optimization problem, which is a class of problems that are hard to solve but easy to verify.

An attempt to address these problems was made by Metropolis et al. in 1953, who developed the Metropolis algorithm, which is a Monte Carlo method that allows to sample configurations of the system according to the Boltzmann distribution. The Metropolis algorithm is a Markov chain Monte Carlo method, which means that it generates a sequence of configurations that are correlated with each other, but the correlation decays exponentially with the number of steps.

It was a revolution in how peaple conceaved mathematics: machine started to be used to solve problems that were not solvable analytically, and the Monte Carlo method was a way to use random numbers to solve problems that were not solvable deterministically.

computing partition function in low temperature limit... organizing principle... probability theory...

\subsection{Sampling a Probability Distribution}

Mean and variance, finite number of samples

good estimator of the mean is the average

a good estimator of the variance is the ``error'', final formula... we can trust the error iff the \(A_i\) are uncorrelated.

\subsection{Hit or Miss: Evaluating the Area of a Circle}

[...]\TODO{draw unit circle inscribed in a square, with random points inside the square, some in the circle and some outside the circle.}

\subsubsection{Estimator for the Average}

[...]

\subsubsection{Confidence Interval for the Average}

computation, check professor and notes

the distribution may not be gaussian, but the central limit theorem guarantees that the distribution of the average is gaussian, and therefore we can use the standard deviation to compute the confidence interval for the average.

independence of observations is eady to achive in laboratory experiments, but in numerical simulations it is not easy to achieve, since the configurations are correlated with each other. The problem of how to generate uncorrelated configurations is a central problem in Monte Carlo simulations, and it is addressed by using Markov chain Monte Carlo methods, which generate a sequence of configurations that are correlated with each other, but the correlation decays exponentially with the number of steps.

\subsection{Direct Sampling}

integral...

probability distribution...

sampled function...

simpson method... error of this methods scales as a power of the number of chunks in which we divide the integration domain, and the exponent depends on the dimensionality...

integrals and sums are basically the same thing in this context, and we need to compute the partition function, which is a sum over an exponential number of configurations.

for the case of the circle...

this was an example of direct or homogeneus sampling, where we sample the whole domain of integration uniformly.

\subsubsection{Failure of homogeneous sampling}

[...]

the choice of the sampling distribution is crucial for the efficiency of the Monte Carlo method, and it is a central problem in Monte Carlo simulations. The choice of the sampling distribution is a trade-off between the ease of sampling and the efficiency of the method. The ideal sampling distribution is the one that is proportional to the integrand, which minimizes the variance of the estimator, but it may be difficult to sample from this distribution.

we can choose a pd which is similar to the integrand, but less explosive, and we can use importance sampling to correct for the difference between the sampling distribution and the integrand...

A smart choice of the sampling distribution is one that makes the sample function as smooth as possible, which minimizes the variance of the estimator. This means we want to transfer as much of the complexity of the integrand to the sampling distribution, which is a trade-off between the ease of sampling and the efficiency of the method.

\subsubsection{Motivation for Markov Chain Monte Carlo}

We can imagine to substitute the distribution with program sampling a sequence of configurations that are correlated with each other, by letting him move to the points in which the function is larger, and therefore the sampling distribution is similar to the integrand. This is the idea behind Markov chain Monte Carlo methods, but the correlation between the configurations is a problem, since it introduces a bias in the estimator.

We need to find a way to select a chain which is representative of the distribution, and we need to find a way to decorrelate the configurations, which is a central problem in Monte Carlo simulations. The problem of how to generate uncorrelated configurations is addressed by using Markov chain Monte Carlo methods, which generate a sequence of configurations that are correlated with each other, but the correlation decays exponentially with the number of steps.

\section{Markov Chain Monte Carlo}

One possibility to sample a probability distribution is to engineer a chain of events spanning the L-dimensional configuration target space: ...

    [...]

We need to find an algorithm that implements these ideas in a practical way, and we need to find a way to decorrelate the configurations, which is a central problem in Monte Carlo simulations.

\subsection{Transition Matrix}

[...]

rewriting the probability as products of transition matrices does not guarantee that each step have memory only of the previous step, since the transition matrix may depend on the whole history of the chain. We need to find a way to construct a transition matrix that depends only on the current configuration, which is a central problem in Markov chain Monte Carlo methods.